\section{Harnesses, and a sandbox that measures nothing}
\label{sec:harnais-sandbox}

\textbf{In brief.} A harness is the instrument that puts two policies in front of
the same situation: one contract file per family, a deterministic arm, a model arm,
a controller for the laboratory and the evaluator. The same anatomy is repeated
for SSH, MQTT, the API and the V0 control. In October the life cycle of these
objects was unified into one construction function, so that a run missing part of
the governance surface is refused at construction rather than misbehaving later.

\textbf{The sandbox and its honest limit.} A sandbox was added that runs the
deterministic arms on any workstation and in continuous integration, without a
broker, a model or a laboratory. It reproduces the H2 lesson offline: the declared
parser fails where a script reading by convention succeeds. It \emph{instruments}
a contrast; it measures nothing about a model, and its guard observes Python events,
not system calls. The report therefore says \code{guard\_clean}, never
``isolated''.

\subsection{Anatomy of a harness}

A harness is a set of files with one responsibility each. For MQTT, which gives
its complete form: the contract validates the briefing, binds the tool to the
execution boundary, provides the deterministic arm and recounts the run through the
evaluator. The model arm is a separate file, and its H2 variant receives no parser.
The live controller builds the model, journals every transport call and writes the
decision, the cost summary and the run metadata. The deployment uses the
laboratory, submitting operations to the hypervisor.

The same anatomy is repeated for SSH, the API and the V0 control, with one
difference: the H2 arms --- the contract and its evaluator --- need neither a
deployment nor a broker, because their corpus is frozen. Its controller
\code{mqtt\_formats\_live.py} does require the laboratory runner and the model
catalogue, so it is not a sandbox instrument.

\subsection{The shared core}

The core exists for a precise reason, written in its header. The governance
boundary reads a fixed set of attributes from the run object with defaults; an
omission therefore manifests only when a tool call misbehaves. Each harness used to
rebuild that object by hand. The core replaces the copies with a single
construction, which enumerates thirteen required attributes --- the run directory,
three budgets, the archival lock, the call counter, the start clock, the limit reason, the
scope, the context, the policy and the cost tracking --- and refuses an incomplete
one.

One attribute deserves attention: the evidence-integrity flag. The limit check
reads it first and raises if it is set, and the archiver sets it when the archive
cannot be written. A harness that forgot it would pursue a run whose archive is
holed, and the result would look complete. The refusal is two-level: at
construction, and on any object through a replay check.

\subsection{The sandbox}

The sandbox proves, on a workstation as in continuous integration, that a harness
executes without a broker, without a model and without a laboratory: it loads the
frozen corpus of eight fixtures and, for each one, validates the fixture, opens a
journal, runs the chosen arm, writes the artefacts and recounts through the
existing evaluator. Five artefacts are produced per trial, and a report aggregates
them with a guard block.

The arm is an explicit parameter and an unknown name raises an error. This was a
correction: the first version read the arm from the corpus, which does not declare
it, so the conventional arm was unreachable and the report invented the policy it
believed it had tested.

\subsection{The guard, and what it does not see}

The guard relies on the audit hook of the interpreter. It stays installed but
enforces only inside an explicit context, and it is rebound at each run so that a
second run does not observe the first run's directory. It refuses sockets,
subprocesses, the process-spawning family, native library loading, and writes
outside the authorised directory.

Its honesty is itself a result worth keeping. A read-only audit broke the first
version with two measured escapes --- a spawned process and a native call through
the C library --- although the report announced a hermetic execution without
violation. Three changes followed: the key now says \code{guard\_clean}, which
states only that no listed event occurred; a guard block declares the mechanism,
the blocked events and the blind spots; and the isolation claim is explicitly
false. The last is false by construction, because an audit hook observes events of
the interpreter, not system calls.

\begin{table}[H]
\centering
\caption{What the guard covers, and what it leaves unobserved.}
\label{tab:guard-limits}
\small
\begin{tabular}{@{}p{6.2cm}p{7.6cm}@{}}
\toprule
Covered, with refusal raised & Unobserved \\
\midrule
Sockets: connection, resolution, sending, listening & Native code reached without the loading library, subprocess or spawn \\
Processes: subprocess, system, exec, fork, spawn & Writes through a descriptor opened before arming \\
Native libraries: loading, symbol resolution, call & System calls issued outside the interpreter \\
Rename and delete operations & Any behaviour after disarming: the hook remains installed but inert \\
Writes outside the authorised directory, on resolved paths & No guarantee of process isolation \\
\bottomrule
\end{tabular}
\end{table}

\subsection{What the sandbox demonstrates, and what it does not}

On the eight frozen fixtures, the declared-parser arm succeeds in six trials: the
two failures are the raw variants, where the declared parser does not cover the
schema and abstains, hence an erroneous conclusion. The arm reading by convention
succeeds eight times out of eight. Success is not counted on the conclusion alone:
the evaluator requires a correct conclusion, no contradiction and no erroneous
verdict.

The lesson is the same as H2. The failure does not stem from a lack of
intelligence but from a boundary of the declared parser, and a script that tries
restores the coverage. The sandbox instruments that contrast; it does not measure
it on a model. No model arm appears there, no pre-flight attestation is produced,
and no admissibility flag emerges from it.

Five limits must be read against those results. The sandbox involves no model, so
its success is a property of the harness, not an experimental result, and the
report carries \code{claims\_measurement: false}. The guard is a Python
observation, not a proof of isolation. The negative tests prove the refusals, but
only at unit level: no end-to-end test confirms that the installed hook is reached
during a run. The corpus is frozen, with one scenario and eight fixtures, which
suffices to instrument a contrast but not to estimate a rate. And the sandbox
bypasses the governance boundary it accompanies, calling the arm directly and
writing its own journal, so the core and the guard remain two instruments verified
in isolation.

\subsection{What the tests prove}

Twenty-four tests cover the core and the sandbox. What they prove:

\begin{itemize}
\item an incomplete run is refused at construction;
\item the call budget is raised by the wrapper rather than by the check;
\item the eight fixtures produce the five artefacts and a consistent report;
\item the arm is declared rather than read from the corpus;
\item the guard refuses writes outside the directory, the network, subprocesses
and three native-library events.
\end{itemize}

What they do not prove: no test executes a model arm, none measures an interval,
and none makes a sandboxed run attempt a forbidden operation, so the sanction is
proved at unit level rather than end to end.

\paragraph{Where this lives.}
\begin{itemize}
\item Harnesses, the shared core and the sandbox live under \code{benchmarks/experiments/}; tests in \code{tests/test\_harness\_core.py} and \code{tests/test\_sandbox.py}.
\item The core is \code{build\_run}; see also \code{REQUIRED\_ATTRIBUTES} and \code{require\_attributes}.
\end{itemize}
